
\documentclass{article}
\usepackage{colortbl}
\usepackage{makecell}
\usepackage{multirow}
\usepackage{supertabular}

\begin{document}

\newcounter{utterance}

\centering \large Interaction Transcript for game `ifeval', experiment `ifeval', episode 240 with Qwen3.5{-}4B.
\vspace{24pt}

{ \footnotesize  \setcounter{utterance}{1}
\setlength{\tabcolsep}{0pt}
\begin{supertabular}{c@{$\;$}|p{.15\linewidth}@{}p{.15\linewidth}p{.15\linewidth}p{.15\linewidth}p{.15\linewidth}p{.15\linewidth}}
    \# & \multicolumn{2}{c}{Player} && \multicolumn{2}{c}{Game Master} \\
    \hline

    \theutterance \stepcounter{utterance}  
    & & & \multicolumn{4}{p{0.6\linewidth}}{
        \cellcolor[rgb]{0.9,0.9,0.9}{
            \makecell[{{p{\linewidth}}}]{
                \texttt{\tiny{[P1$\langle$GM]}}
                \texttt{Write a casual blog post about how the outer solar system is different from the inner solar system, and what that means for the possibility of life. Wrap your entire response with double quotation marks. Your response should contain 17 or more sentences.} \\
            }
        }
    }
     \\ \\

    \theutterance \stepcounter{utterance}  
    & \multicolumn{4}{p{0.6\linewidth}}{
        \cellcolor[rgb]{0.561,0.737,0.561}{
            \makecell[{{p{\linewidth}}}]{
                \texttt{\tiny{[P1$\rangle$GM]}}
                \texttt{"Have you ever looked up at the night sky and wondered why the planets seem so different from one another? The inner solar system is a hot, rocky playground dominated by Mercury, Venus, Earth, and Mars. These worlds are close to the sun, which means they are scorching and dry with little to no atmosphere to keep them cool. In contrast, the outer solar system is a frozen, icy frontier filled with gas giants and distant dwarf planets. Jupiter, Saturn, Uranus, and Neptune are massive balls of hydrogen and helium that spin incredibly fast. The temperatures down there are so low that water exists as ice rather than liquid or steam. This huge difference in temperature creates a completely different environment for any potential life forms. Life as we know it needs liquid water to survive, and that is mostly found in the middle of our solar system. The inner planets have solid surfaces that you could theoretically walk on if you had a spacesuit. However, the outer planets are mostly gas and liquid, making them impossible to land on with current technology. Some scientists believe that life might exist in the subsurface oceans of moons like Europa or Enceladus. These icy moons are actually orbiting the gas giants, not the sun directly, which is a fascinating twist. They are heated by the gravitational pull of their large neighbors, keeping their oceans liquid deep below the crust. This hidden warmth could be the key to finding alien microbes hiding in the dark and cold. While the inner system offers rocky terrain, the outer system offers vast, liquid worlds waiting to be explored. Both regions hold secrets that humanity is only just beginning to uncover with our new telescopes. The search for life is not limited to one type of planet but spans the entire diversity of our cosmic neighborhood. It is truly exciting to think about what strange creatures might be hiding in the icy depths of the outer reaches. We are on the verge of a new era where we might finally find our neighbors in the stars."} \\
            }
        }
    }
    & & \\ \\

    \theutterance \stepcounter{utterance}  
    & & & \multicolumn{2}{p{0.3\linewidth}}{
        \cellcolor[rgb]{0.9,0.9,0.9}{
            \makecell[{{p{\linewidth}}}]{
                \texttt{\tiny{[GM$|$GM]}}
                {'startend:quotation': {}, 'length_constraints:number_sentences': {'relation': 'at least', 'num_sentences': 17}}
            }
        }
    }
    & & \\ \\

    \theutterance \stepcounter{utterance}  
    & & & \multicolumn{2}{p{0.3\linewidth}}{
        \cellcolor[rgb]{0.9,0.9,0.9}{
            \makecell[{{p{\linewidth}}}]{
                \texttt{\tiny{[GM$|$GM]}}
                \texttt{game\_result = WIN} \\
            }
        }
    }
    & & \\ \\

\end{supertabular}
}

\end{document}
